\section{Three connected research questions}
In a 20-minute research-brief session, an Evidence Agent flags factual or citation issues, a Planning Agent flags deadline or sequencing risks, and an Idea Agent proposes writing options. They submit notifications but cannot interrupt directly. An Attention Coordinator must interrupt now, queue a message for the next five-minute checkpoint, or defer it to the minute-20 digest. The integrated question is how to combine current importance and urgency claims with task-specific historical reliability while preserving human oversight.

\textbf{Q1---Economics:} When does historical reliability improve allocation of scarce attention, and who bears the costs of false alarms, missed warnings, and interruptions?\par
\textbf{Q2---Computation:} Can a transparent, reproducible history-calibrated scheduler reduce unnecessary interruptions without delaying too many important messages?\par
\textbf{Q3---Behavior:} When users respond less to an agent, is the cause cognitive burden from repeated interruption or declining trust after repeated false alarms?

The scheduling rule allocates attention, changes agents' incentives and voice, and shapes burden and trust. Figure~\ref{fig:teaser} summarizes how current claims, historical reliability, and the three disciplinary questions connect to interruption authority.

\section{Answer to the economics question}
This proposed allocation problem draws intellectual lineage from incomplete-information games~\cite{harsanyi1968}; this PS1 does not solve a Nash or Bayesian equilibrium. The user, agents, and coordinator are players. Reports are actions; an agent's reporting policy is its strategy. The coordinator sees claims and task-matched history, not true message value. Welfare balances timely valid warnings against interruption, delay, and missed-warning costs. A future extension may model incentives to inflate claims, but the current heuristic neither solves nor demonstrates strategic misreporting.

Normalize importance $I$, urgency $U$, and reliability $R$ to $[0,1]$. The baseline is $N=0.5I+0.5U$ and $S=N$; the proposed score is $S=NR$. Multiplication is neither novel nor presumed optimal. The falsifiable hypothesis is that task-specific reliability helps when past validity predicts current validity, but harms welfare with sparse histories, domain shifts, or discounted rare-warning agents. Audits must report whose messages are delayed, not only attention saved. Nash~\cite{nash1950} provides intellectual lineage and a possible future benchmark, not a solved result here.

\section{Answer to the computational question}
An authored deterministic diagnostic evaluates 18 synthetic notifications. Online routing uses current claims and Beta-smoothed task-specific history; hidden validity and loss labels are evaluation-only. Shared thresholds select interrupt, next checkpoint, or digest.

Current claims produced 10 immediate interruptions, timely recall of all eight important alerts, and joint loss 21.15. Multiplication reduced interruptions to 2 and attention cost from 21.15 to 5.40, but delayed alerts lowered recall to 25\% and raised joint loss to 54.40. It is therefore not better in this diagnostic: it saved attention by suppressing messages enough to delay important warnings.

With ten immediate slots and identical queues, both rankings had joint loss 21.60; Idea received four slots under current claims and zero under reputation. These synthetic results establish neither generality, optimality, real-agent behavior, nor deployed safety.

\clearpage\balance
\section{Answer to the behavioral question}
\emph{Burden} predicts disengagement after frequent valid interruptions; \emph{trust loss} predicts disengagement after false alarms at a fixed interruption rate. A future consented study could vary frequency while holding validity constant, then vary false-alarm history while holding frequency constant. Response, recall, completion, and perceived trust could distinguish them; no human-subject evidence exists.

If burden dominates, batching or interruption budgets may outperform reputation. If trust loss dominates, explanations and reputation recovery matter more. Equal responses to reliable and unreliable agents would weaken the case for $R$. Tests must also detect self-reinforcing silence: fewer interruptions create fewer chances to demonstrate validity.

\section{Advanced development and future directions}
Nobel laureates John Nash~\cite{nash1950} and John Harsanyi~\cite{harsanyi1968} clarify strategic interdependence and hidden information; an official award-record citation remains pending. They do not determine how specialized AI agents earn access to scarce human oversight. The contribution is a falsifiable framework joining allocation incentives, an auditable scheduler, and burden-versus-trust mechanisms. A Turing Award lineage remains pending source verification.

\begin{table}[h]\caption{From laureate foundations to a new interdisciplinary contribution.}
\begin{tabularx}{\columnwidth}{@{}p{1.6cm}Y@{}}\toprule
\textbf{Horizon}&\textbf{Milestone and evidence}\\\midrule
Foundation&Nash and Harsanyi: strategic choice under interdependence and incomplete information.\\
PS1&Specify $S=N$ versus $S=NR$ and audit a reproducible synthetic 20-minute scenario.\\
Next study&Compare schedulers and separate interruption burden from trust loss.\\
Longer term&Test domain shifts, reputation recovery, oversight, and unequal voice.\\
2056&Establish a falsifiable framework for accountable interruption authority.\\\bottomrule
\end{tabularx}\end{table}

\noindent\textbf{Expected 2056 Nobel Prize and Turing Award contribution statement.}
By 2056, I aspire to contribute a falsifiable framework for allocating interruption authority among AI agents. I will integrate incentive analysis, transparent scheduling, and behavioral evidence about burden and trust. Progress requires reproducibility, cross-domain evidence, reputation recovery, and audits of missed warnings and unequal voice. This ambition neither declares $S=NR$ optimal nor predicts an award.

\subsection*{Open Science Statement}
\textbf{GitHub:} \GitHubLink\\
\textbf{Google Colab:} \ColabLink\par
The repository contains scheduler code, authored synthetic inputs, tests, saved outputs, and an executable Colab notebook. A fresh Colab Run All check succeeded. These materials reproduce the constructed diagnostic; they are not human or real-agent evidence.

\subsection*{Statement of Contribution to the UN's SDGs}
\textbf{Hugging Face:}\\
\DemoLink\\
The deployed open-access simulator supports exploration of current-claim-only and history-calibrated attention rules in an interactive scenario relevant to \textbf{SDG 4---Quality Education}~\cite{sdg4}. It provides an accessible diagnostic interface, but no measured educational benefit is claimed.

\subsection*{Submission milestones}
First draft: September 13, 2026, 11:00 P.M. Beijing time (UTC+8). Class review: September 14. Proposed: rebuttal September 16 and revised submission September 20, both at 11:00 P.M. Update Appendix~\ref{app:abstract} with each reflection and revision.
